% 原分章PDF115页附录7.J标题—118页末，印刷页691—694。
\section*{附录7.J\quad 贝叶斯线性模型预测分布的推导}
\addcontentsline{toc}{section}{附录7.J 贝叶斯线性模型预测分布的推导}
\markright{附录7.J 贝叶斯线性模型预测分布的推导}

为使叙述自洽，本附录按照Rasmussen and Williams（2016）的方法，简要推导贝叶斯线性模型中的预测分布。由式（7.367）、（7.373），得到
\begin{equation}
\begin{aligned}
&p(y_{new}\mid x_{new},w)p(w\mid X,y)\\
&\quad\propto\exp\biggl\{-\frac12\biggl[
 \frac1{\sigma_n^2}(y_{new}-x_{new}^Tw)^T(y_{new}-x_{new}^Tw)
 +(w-\mu)^T\Lambda(w-\mu)\biggr]\biggr\}.
\end{aligned}\tag{7.J.1}
\end{equation}
注意
\begin{equation}
\begin{aligned}
&\frac1{\sigma_n^2}(y_{new}-x_{new}^Tw)^T(y_{new}-x_{new}^Tw)\\
&\quad=\frac1{\sigma_n^2}y_{new}^2-\frac2{\sigma_n^2}w^Tx_{new}y_{new}
 +\frac{w^Tx_{new}x_{new}^T}{\sigma_n^2}w,
\end{aligned}\tag{7.J.2}
\end{equation}
\begin{equation}
(w-\mu)^T\Lambda(w-\mu)=w^T\Lambda w-2w^T\mu+\mu^T\Lambda\mu.\tag{7.J.3}
\end{equation}
将式（7.J.2）、（7.J.3）代入式（7.J.1），得到
\begin{equation}
\begin{aligned}
&p(y_{new}\mid x_{new},w)p(w\mid X,y)\\
&\quad\propto\exp\biggl\{-\frac12\biggl[
 w^T\left(\Lambda+\frac{x_{new}x_{new}^T}{\sigma_n^2}\right)w\\
&\hspace{8em}{}-2w^T\left(\Lambda\mu+\frac{x_{new}y_{new}}{\sigma_n^2}\right)
 +\frac{y_{new}^2}{\sigma_n^2}\biggr]\biggr\}.
\end{aligned}\tag{7.J.4}
\end{equation}
令
\begin{equation}
V=\Lambda+\frac{x_{new}x_{new}^T}{\sigma_n^2}.\tag{7.J.5}
\end{equation}
于是，
\begin{equation}
V^{-1}=\Lambda^{-1}-\frac1{\sigma_n^2}
 \frac{\Lambda^{-1}x_{new}x_{new}^T\Lambda^{-1}}
 {1+\dfrac1{\sigma_n^2}x_{new}^T\Lambda^{-1}x_{new}}.\tag{7.J.6}
\end{equation}
定义
\begin{equation}
(w-m)^TV(w-m)=w^TVw-2w^TVm+m^TVm.\tag{7.J.7}
\end{equation}
为求均值 $m$，令
\begin{equation}
Vm=\Lambda\mu+\frac{x_{new}y_{new}}{\sigma_n^2}.\tag{7.J.8}
\end{equation}
由式（7.J.8）解出 $m$，得到
\begin{equation}
m=V^{-1}\left(\Lambda\mu+\frac{x_{new}y_{new}}{\sigma_n^2}\right),\tag{7.J.9}
\end{equation}
从而
\begin{equation}
\begin{aligned}
m^TVm
 &=\left(\Lambda\mu+\frac{x_{new}y_{new}}{\sigma_n^2}\right)^T
 V^{-1}\left(\Lambda\mu+\frac{x_{new}y_{new}}{\sigma_n^2}\right)\\
 &=\frac{x_{new}^TV^{-1}x_{new}}{\sigma_n^4}y_{new}^2
 +2\frac{\mu^T\Lambda V^{-1}x_{new}y_{new}}{\sigma_n^2}+\mathrm{const}.
\end{aligned}\tag{7.J.10}
\end{equation}
将式（7.J.10）代入式（7.J.4），得到
\begin{equation}
\begin{aligned}
&p(y_{new}\mid x_{new},w)p(w\mid X,y)\\
&\quad\propto\exp\left\{-\frac12\left[(w-m)^TV(w-m)-m^TVm
 +\frac{y_{new}^2}{\sigma_n^2}\right]\right\}\\
&\quad\propto\exp\biggl\{-\frac12(w-m)^TV(w-m)\\
&\hspace{6em}{}-\frac12\biggl[\frac{y_{new}^2}{\sigma_n^2}
 -\frac{x_{new}^TV^{-1}x_{new}}{\sigma_n^4}y_{new}^2\\
&\hspace{11em}{}-2\frac{\mu^T\Lambda V^{-1}x_{new}y_{new}}{\sigma_n^2}
 +\mathrm{const}.\biggr]\biggr\}.
\end{aligned}\tag{7.J.11}
\end{equation}
将式（7.J.11）代入式（7.375），得到
\begin{equation}
\begin{aligned}
&p(y_{new}\mid x_{new},X,y)
 =\int p(y_{new}\mid x_{new},w)p(w\mid X,y)\,dw\\
&\quad\propto\exp\biggl\{-\frac12\biggl[\frac1{\sigma_n^2}
 \left(1-\frac1{\sigma_n^2}x_{new}^TV^{-1}x_{new}\right)y_{new}^2\\
&\hspace{12em}{}-2\frac{\mu^T\Lambda V^{-1}x_{new}y_{new}}{\sigma_n^2}
 \biggr]\biggr\}.
\end{aligned}\tag{7.J.12}
\end{equation}
令
\[
\begin{aligned}
\lambda(y_{new}-\bar y)^2
={}&\frac1{\sigma_n^2}\left(1-\frac1{\sigma_n^2}x_{new}^TV^{-1}x_{new}\right)y_{new}^2\\
&{}-2\frac{\mu^T\Lambda V^{-1}x_{new}y_{new}}{\sigma_n^2}+\bar y^2,
\end{aligned}
\]
得到
\begin{equation}
\lambda=\frac1{\sigma_n^2}\left(1-\frac1{\sigma_n^2}x_{new}^TV^{-1}x_{new}\right)\tag{7.J.13}
\end{equation}
以及
\begin{equation}
\bar y=\frac1\lambda\frac1{\sigma_n^2}x_{new}^TV^{-1}\Lambda\mu,\tag{7.J.14}
\end{equation}
其中
\[
\mu=\frac1{\sigma_n^2}\Lambda^{-1}Xy.
\]
注意
\begin{equation}
\begin{aligned}
x_{new}^TV^{-1}x_{new}
={}&x_{new}^T\Lambda^{-1}x_{new}
 -\frac1{\sigma_n^2}\frac{(x_{new}^T\Lambda^{-1}x_{new})^2}
 {1+\dfrac1{\sigma_n^2}x_{new}^T\Lambda^{-1}x_{new}}\\
={}&x_{new}^T\Lambda^{-1}x_{new}\left[1-\frac1{\sigma_n^2}
 \frac{x_{new}^T\Lambda^{-1}x_{new}}
 {1+\dfrac1{\sigma_n^2}x_{new}^T\Lambda^{-1}x_{new}}\right]\\
={}&\frac{x_{new}^T\Lambda^{-1}x_{new}}
 {1+\dfrac1{\sigma_n^2}x_{new}^T\Lambda^{-1}x_{new}}.
\end{aligned}\tag{7.J.15}
\end{equation}
将式（7.J.15）代入式（7.J.13），得到
\begin{equation}
\lambda=\frac1{\sigma_n^2\left(1+\dfrac1{\sigma_n^2}x_{new}^T\Lambda^{-1}x_{new}\right)}\tag{7.J.16}
\end{equation}
或
\begin{equation}
\frac1\lambda=\sigma_n^2\left(1+\frac1{\sigma_n^2}x_{new}^T\Lambda^{-1}x_{new}\right).\tag{7.J.17}
\end{equation}
回顾
\begin{equation}
\begin{aligned}
\bar y&=\frac1\lambda\frac1{\sigma_n^2}x_{new}^TV^{-1}\Lambda\mu\\
 &=\mu^T\frac{\Lambda V^{-1}x_{new}}{\lambda^T\lambda\sigma_n^2}.
\end{aligned}\tag{7.J.18}
\end{equation}
接下来证明
\[
\frac{\Lambda V^{-1}x_{new}}{\lambda\sigma_n^2}=x_{new}.
\]
为此，作如下变形：
\begin{equation}
\begin{aligned}
V\Lambda^{-1}x_{new}
 &=\left(\Lambda+\frac{x_{new}x_{new}^T}{\sigma_n^2}\right)\Lambda^{-1}x_{new}\\
 &=x_{new}+\frac{x_{new}x_{new}^T\Lambda^{-1}x_{new}}{\sigma_n^2}\\
 &=\left(1+\frac{x_{new}^T\Lambda^{-1}x_{new}}{\sigma_n^2}\right)x_{new}.
\end{aligned}\tag{7.J.19}
\end{equation}
由式（7.J.17），有
\begin{equation}
\frac1{\lambda\sigma_n^2}
 =\frac1{\lambda\sigma_n^2}x_{new}^T\Lambda^{-1}x_{new}.\tag{7.J.20}
\end{equation}
将式（7.J.20）代入式（7.J.19），得到
\begin{equation}
V\Lambda^{-1}x_{new}=\frac1{\lambda\sigma_n^2}x_{new}.\tag{7.J.21}
\end{equation}
将式（7.J.21）两边乘以 $\Lambda V^{-1}$，得到
\begin{equation}
x_{new}=\frac{\Lambda V^{-1}x_{new}}{\lambda\sigma_n^2}\tag{7.3A22}
\end{equation}
将式（7.J.22）代入式（7.J.18），得到
\begin{equation}
\begin{aligned}
\bar y&=\mu^Tx_{new}\\
 &=x_{new}^T\mu.
\end{aligned}\tag{7.J.23}
\end{equation}
结合式（7.J.14）、（7.J.23），得到
\begin{equation}
\bar y=x_{new}^T\frac1{\sigma_n^2}\Lambda^{-1}Xy.\tag{7.J.24}
\end{equation}
\begin{equation}
\begin{aligned}
p(y_{new}\mid x_{new},X,y)
 &\propto\exp\left\{-\frac{\left(y_{new}-\dfrac1{\sigma_n^2}x_{new}^T\Lambda^{-1}Xy\right)}
 {2x_{new}^T\Lambda^{-1}x_{new}}\right\}\\
 &=N\left(\frac1{\sigma_n^2}x_{new}^T\Lambda^{-1}Xy,\Lambda^{-1}x_{new}\right).
\end{aligned}\tag{7.J.25}
\end{equation}
